\section{问题分析}

\subsection{总体思路}

本题四问形成一条完整的推理链：\textbf{问题一}把``数据质量''这一抽象概念量化为
可计算的评分体系，并建立配比与损失之间的数量关系；\textbf{问题二}把质量维度
嵌入经典标度律，得到可微的广义能力方程 $L(N,D,Q)$；\textbf{问题三}以该方程为
目标函数，在算力约束下求解最优资源配置并刻画结构性转移；\textbf{问题四}跳出
单模型视角，站在行业层面分解技术演进并把问题一--三得到的``能力--资源''关系
用于前沿预测。四问之间通过``质量分 $Q$''与``损失 $L$''两个量纲紧密衔接。

\subsection{问题一的分析}

质量评分的关键难点有三：指标量纲与方向不统一、指标之间相互冲突（如``内容
相关''与``格式干净''在同一文档上评分相反）、以及配比效应随模型规模转移。
对策上分别采用\textbf{组合赋权 + TOPSIS}（消除主观性与量纲差异）、
\textbf{双族冲突指数 + 对称截尾融合}（显式刻画并消解冲突）、以及
\textbf{跨尺度系数收缩分析}（把配比--损失关系放在多规模数据上检验其系数随
规模的幂律衰减）。域级质量按样本权重聚合，并用 1M 抽样集与全量扩展集对照
验证评分的可迁移性。

\subsection{问题二的分析}

经典标度律 $L=E+A N^{-\alpha}+B D^{-\beta}$ 只含规模变量，质量 $Q$ 需以
可退化的方式加入：当 $Q=1$（理想质量）时应退化为经典形式，因此质量项取
$(1-Q)^{g}$ 结构。在数据侧，B1 轨迹拟合经典参数，B4/B5 做跨家族验证
（注意各家族验证集与词表不同造成的 Loss 尺度偏移，需用家族偏移修正 $R^2$），
B6/B7 拟合广义形式，B8 作为外推稳健性对照（须先诊断其 $Q$ 方向与 B6/B7
是否一致）。弹性与等价性分析（如``质量提升 0.1 等价于多少参数''）为问题三
的资源配置提供直觉。

\subsection{问题三的分析}

在预算约束下，$N$、$D$、$Q$ 通过成本方程耦合：$6ND$ 为训练成本，
$D[g(Q)-g(Q_0)]_+$ 为质量提升成本（只有超过基线 $Q_0$ 才产生），
$\eta N D L_{\text{ctx}}$ 为注意力成本。后两项的引入使优化问题从经典的
``$ND$ 乘积约束''推广为非线性约束。结构性转移的数学定义采用
\textbf{成本份额主导成分切换}与\textbf{质量变量是否离开边界 $Q_0$}两类指标：
当预算连续增长时，最优解从``纯规模扩张（$Q=Q_0$）''进入``质量投入激活
（$Q>Q_0$）''，再进入``质量饱和（$Q=1$）''的阶段性变化，转移点由数值扫描
配合 KKT 条件识别。

\subsection{问题四的分析}

前沿能力 $S$ 与规模、时间的关系用 90\% 分位数回归 $\ln S=c+b_N\ln N+b_T t$
刻画：$b_N\ln N$ 项对应规模贡献，$b_T t$ 项对应非规模技术进步（算法、对齐、
工程与数据组织等）。预测时先由 C4 估计开源模型参数量增速 $g_N$，再按
``基础情景''（增速延续）与``放缓情景''（增速减半）外推，用自助法给出 90\%
置信区间。Loss--Benchmark 桥接（C6）把问题二得到的损失预测映射回排行榜分数，
实现``损失--能力''两套度量的统一。
